\chapter{Solving Quadratics}

A quadratic function has three terms: $ax^2 + bx + c$. $a$, $b$, and
$c$ are known as the \textit{coefficients}. The coefficients can be any
constant, except that $a$ can never be zero. (If $a$ is zero, it is a linear
function, not a quadratic.)

When you have an equation with a quadratic function on one side and a
zero on the other, you have a quadratic equation. For example:

$$72x^2 - 12x + 1.2 = 0$$

How can you find the values of $x$ that will make this equation true?

You can always reduce a quadratic equation so that the first
coefficient is 1, so that your equation looks like this:

$$x^2 +bx + c = 0$$

For example, if you are asked to solve $4x^2 + 8x - 19 = -2x^2 - 7$
\begin{multline*}
  4x^2 + 8x - 19 = -2x^2 - 7 \\
  6x^2 + 8x -12 = 0 \\
  x^2 + \frac{4}{3}x - 2 = 0
\end{multline*}
Here, $b = \frac{4}{3}$ and $c = -2$.
% This looks a little wonky, maybe add explanations too

\begin{mdframed}[style=important]
$x^2 + bx + c = 0$ when
\begin{equation*}
x = -\frac{b}{2} \pm \frac{\sqrt{b^2 - 4c}}{2}  
\end{equation*}
Only when $a=1$
\end{mdframed}

What does this mean?

For any $b$ and $c$, the graph of $x^2 + bx + c$ is a parabola
that goes up on each end. Its low point, the vertex, is at $x = -\frac{b}{2}$.
For a quadratic with any value of $a$, the vertex is at $x = -\frac{b}{2a}$,
which is called the vertex formula\index{vertex formula}.

If there are no real roots ($b^2 - 4c < 0$), which means the
parabola never gets low enough to cross the $x$-axis:
\begin{center}
  \begin{tikzpicture}
      \begin{axis}[
          xmin=-0.5,xmax=2.75,
          ymin=-1,ymax=5,
          axis x line=middle,
          axis y line=middle,
          axis line style=<->,
          xlabel={$x$},
          ylabel={$y$},
        ]
        \addplot[no marks,sdkblue] expression[domain=-0.25:2.5,samples=100]{x^2 - 2* x + 3} node[above, xshift=-1cm] {$x^2 - 2x + 3$};
        \addplot[dashed,gray] coordinates {(1,-1)(1,3)};
      \end{axis}
  \end{tikzpicture}
\end{center}

If there is one real root ($b^2 - 4c = 0$), it means that the parabola only touches the x-axis.
\begin{center}
  
  \begin{tikzpicture}
      \begin{axis}[
          xmin=-0.5,xmax=3.75,
          ymin=-0.5,ymax=5,
          axis x line=middle,
          axis y line=middle,
          axis line style=<->,
          xlabel={$x$},
          ylabel={$y$},
        ]
        \addplot[no marks,sdkblue] expression[domain=-0.25:3.5,samples=100]{x^2 - 4*x + 4} node[above, xshift=-1cm] {$x^2 - 4x + 4$};
        \addplot[dashed,gray] coordinates {(2,-0.5)(2,1)};
      \end{axis}
  \end{tikzpicture}
\end{center}

If there are two real roots ($b^2 - 4c > 0$), it means that the parabola crosses the x-axis twice as it dips below and then returns:
\begin{center}
  
  \begin{tikzpicture}
    \begin{axis}[
      xmin=-0.5,xmax=3.75,
      ymin=-2,ymax=5,
      axis x line=middle,
      axis y line=middle,
      axis line style=<->,
      xlabel={$x$},
      ylabel={$y$},
      ]
      \addplot[no marks,sdkblue] expression[domain=-0.25:3.5,samples=100]{x^2 - 3*x + 1} node[above, xshift=-1cm] {$x^2 - 3x + 1$};
      \addplot[dashed,gray] coordinates {(1.5,-2)(1.5,1)};
    \end{axis}
  \end{tikzpicture}
  
\end{center}
\begin{Exercise}[title={Roots of a Quadratic}, label=solve_quadratic]

  In the last chapter, you found that the function for the height of your flying hammer is:

  $$p = -\frac{1}{2}9.8 t^2 + 12t + 2$$

  At what time will the hammer hit the ground?

  
\end{Exercise}
\begin{Answer}[ref=solve_quadratic]

  For what $t$ is  $-4.9 t^2 + 12t + 2 = 0$?  Start by dividing both sides of the equation by -4.9.

  $$t^2 - 2.45 t - 0.408 = 0$$

  The roots of this are at

  $$x = -\frac{b}{2} \pm \frac{\sqrt{b^2 - 4c}}{2} = -\frac{-2.45}{2} \pm \frac{\sqrt{(-2.45)^2 - 4(-0.408)}}{2} = 1.225 \pm 1.382$$

  We only care about the root after we release the hammer ($t > 0$).

  $1.225 + 1.382 \approx 2.61$ seconds after releasing the hammer, it will hit the ground.

  
\end{Answer}


\section{The Traditional Quadratic Formula}

If the last explanation was a little tricky to understand, the quadratic formula is a nifty tool.

\begin{mdframed}[style=important, frametitle={The Quadratic Formula}]

$ax^2 + bx + c = 0$ when
\begin{equation*}
  x = \frac{-b \pm \sqrt{b^2 - 4ac}}{2a}
\end{equation*}

\end{mdframed}

This formula works for any quadratic equation, including ones where $a$ is
not 1. Recall the equation from the start of this chapter,
$6x^2 + 8x - 12 = 0$. Here $a = 6$, $b = 8$, and $c = -12$. Substitute these
into the formula:
\begin{equation*}
  x = \frac{-8 \pm \sqrt{8^2 - 4(6)(-12)}}{2(6)}
    = \frac{-8 \pm \sqrt{64 + 288}}{12}
    = \frac{-8 \pm \sqrt{352}}{12}
\end{equation*}
Since $\sqrt{352} \approx 18.762$, the $+$ gives
$x \approx \frac{10.762}{12} \approx 0.897$, and the $-$ gives
$x \approx \frac{-26.762}{12} \approx -2.230$.

Watch the signs. Since $c = -12$, the term $-4ac$ becomes
$-4(6)(-12) = +288$, so the number under the square root grows.

Earlier, you divided this equation by 6 to get $x^2 + \frac{4}{3}x - 2 = 0$.
The formula for $a = 1$ gives the same roots: with $b = \frac{4}{3}$ and
$c = -2$, $-\frac{b}{2} \pm \frac{\sqrt{b^2 - 4c}}{2}$ also comes out to
$x \approx 0.897$ and $x \approx -2.230$. Use whichever form is easier for
the numbers in front of you.

The formula also answers the question at the start of this chapter. For
$72x^2 - 12x + 1.2 = 0$, the number under the square root is
$(-12)^2 - 4(72)(1.2) = 144 - 345.6 = -201.6$. No real number is the square
root of a negative number, so this equation has no real solutions.


\section{Factoring Quadratics}

The quadratic formula always works, but the arithmetic can get heavy,
especially without a calculator. For some quadratics, factoring is faster.

Factoring rewrites a quadratic as a product of two simpler expressions. For
example, $x^2 + 5x + 6$ can be written as $(x + 2)(x + 3)$. Expanding the
product gives back the original: $x^2 + 3x + 2x + 6 = x^2 + 5x + 6$.

The factored form makes the roots easy to read because of the
\newterm{zero product property}\index{zero product property}: if two numbers
multiply to zero, at least one of them must be zero. So
$(x + 2)(x + 3) = 0$ means $x + 2 = 0$ or $x + 3 = 0$, which gives $x = -2$
or $x = -3$. These are the same roots the quadratic formula gives, found
without a square root.

Most quadratics you factor by hand have $a = 1$, in the form
$x^2 + bx + c$. Certain quadratics with $a \neq 1$ also factor, using the
AC method below.

\subsection*{Types of Factoring Techniques}

\begin{enumerate}
  \item Factoring out the GCF (Greatest Common Factor) \\
  Always check for a common factor first:
  \[
  2x^2 + 4x = 2x(x + 2)
  \]

  \item Factoring Trinomials: \( a = 1 \) \\
  For expressions like \( x^2 + bx + c \), find two numbers that:
  \begin{itemize}
    \item Multiply to \( c \)
    \item Add to \( b \)
  \end{itemize}
  \[
  x^2 + 5x + 6 = (x + 2)(x + 3)
  \]

  \item Factoring Trinomials: \( a \neq 1 \) \\
  Use the ``AC method'':
  \begin{enumerate}
    \item Multiply \( a \times c \)
    \item Find two numbers that multiply to \( ac \) and add to \( b \)
    \item Split the middle term and factor by grouping
  \end{enumerate}
  \[
  6x^2 + 11x + 4 = 6x^2 + 3x + 8x + 4 = 3x(2x + 1) + 4(2x + 1) = (3x + 4)(2x + 1)
  \]

  \item Special Cases:
  \begin{itemize}
    \item Perfect Square Trinomials: \\
    \( a^2 + 2ab + b^2 = (a + b)^2 \) \\
    \( x^2 + 6x + 9 = (x + 3)^2 \)

    \item Difference of Squares: \\
    \( a^2 - b^2 = (a - b)(a + b) \) \\
    \( x^2 - 16 = (x - 4)(x + 4) \)
  \end{itemize}
\end{enumerate}

\subsection*{Tips}
\begin{itemize}
  \item Always factor out the GCF first.
  \item Check your result by expanding (FOIL). The inverse of factoring is expanding using the FOIL method.
  \item Some quadratics cannot be factored using integers. For those, use the quadratic formula.
\end{itemize}


%FIXME completing the square section? 

% we should have all 3 factoring methods at least mentioned or taught